% ISMIR 2026 LBD: at most 2 scientific pages + 1 references/acknowledgments/AI page.
\documentclass{article}
\pdfoutput=1
\usepackage[T1]{fontenc}
\usepackage[utf8]{inputenc}
\usepackage[lbd,submission]{ismir}
\usepackage{amsmath,cite,url,graphicx,booktabs,microtype}
\emergencystretch=1em
\def\lbd{}
\newcommand{\jams}{\texttt{jams}}
\title{jams: Inspecting a Modular Audio-to-MIDI Pipeline}
\oneauthor{John Hurliman}{Independent researcher \\ {\tt jhurliman@jhurliman.org}}
\def\authorname{J. Hurliman}
\begin{document}
\maketitle

\begin{abstract}
\jams{} is an open-source local music-analysis service with a browser interface for
inspecting separated audio, beat grids, functional segments, and per-group MIDI.
The demonstration focuses on a practical problem: a cascade can produce plausible
stems while making note errors that matter during editing. Users can replay the
mixture and inspect transcribed notes in aligned piano rolls; separated audio
files are available from the service for external inspection.
An accompanying retrospective Slakh2100 case study compares two transcribers before
and after a fixed separator. Its reported accompaniment F1 favors YourMT3+ over
basic-pitch in both conditions, but does not establish an advantage over direct
mixture transcription. We distinguish this preliminary evidence from the
prototype's interface capabilities and identify the archived outputs needed for
independent recomputation.
\end{abstract}

\section{Purpose and scope}
A modular transcription workflow can provide both editable audio groups and MIDI,
allowing a musician to choose which representation to use. Source separation followed
by transcription is an established approach~\cite{derby2024pipeline}. The difficult
part for an application is making the resulting errors inspectable: interference,
missed attacks, extra notes, and merged instruments can affect the symbolic output
without being obvious from a single aggregate quality score.

\jams{} combines a local HTTP analysis service and a browser-based annotator.\footnote{\raggedright
\url{https://github.com/jhurliman/jams} (distinct from the annotation library
\url{https://github.com/marl/jams}).}
The application targets DJ and music-production workflows, but the numerical
comparison discussed here uses synthetic instrumental audio. We do not infer
performance on real electronic dance music from that dataset. The contribution of
the demo is an integrated inspection workflow and a concrete component-comparison
question, not a new separation or transcription architecture.

\section{System and demonstration}
The service returns global key and tempo estimates and optionally beat/downbeat
times, functional segments, separated audio, and MIDI. The separation path uses
an SCNet variant~\cite{tong2024scnet} with drums, bass, other, and vocals outputs.
The other group combines non-bass pitched instruments; it is not an instrument
label. YourMT3+~\cite{chang2024yourmt3} provides pitched-note transcription, while
a separate local model handles drums. The worker architecture keeps incompatible
model environments in separate processes. Initial setup downloads some model
artifacts; the application should be installed and warmed up before a demonstration.

The demonstration follows an audio example through three operations. First, the
user imports a recording and opens its waveform with the returned beat grid and
segment boundaries. Second, the user replays the mixture while inspecting the MIDI
piano rolls for missing or extra attacks and pitch errors. Separated audio can
be examined in an external player; the browser has no stem-solo mixer. Third, the user adjusts segment boundaries or labels and exports
per-group MIDI for further work in a sequencer. These operations expose the output
for human inspection; the application does not establish transcription correctness
simply by displaying it.

The interface also provides a way to examine an evaluation example with known
references. A dense accompaniment passage is useful for discussing what an
instrument-agnostic score counts: a correctly timed note can be matched even when
instrument identity has been discarded. A bass example illustrates why pitch
postprocessing must be disclosed. The project's fixed octave shift is an empirical
pipeline choice that still requires an audio/MIDI alignment audit, not a universal
property of MIDI bass notation.

\section{Preliminary component comparison}
The project ledger records experiments on the 151-track Slakh2100-redux test
split~\cite{manilow2019slakh}. For each track, non-bass pitched reference stems are
summed into other and their MIDI notes are merged. The main metric is per-track
onset-and-pitch F1, macro-averaged across tracks, with one-to-one matching at 50\,ms
and 50 cents, ignoring note offsets~\cite{raffel2014mireval}. Beat quantization is
disabled. Thus the measure does not assess durations, instrument assignments,
velocities, or editing effort.

\begin{table}[t]
\centering\small
\begin{tabular}{lrr}
\toprule
Other input & basic-pitch & YourMT3+ \\
\midrule
Reference group & 0.4897 & 0.8488 \\
SCNet output & 0.4733 & 0.7877 \\
\bottomrule
\end{tabular}
\caption{Reported mean onset-and-pitch F1 for accompaniment, $n=151$.
These are historical summaries, not newly reproduced measurements.}
\label{tab:comparison}
\end{table}

Replacing basic-pitch~\cite{bittner2022basicpitch} with YourMT3+ raises the reported
mean by 0.3591 on reference input and 0.3144 after the same separator
(Table~\ref{tab:comparison}). The archived reference-input paired 95\% bootstrap
interval is [0.3471, 0.3713]; no paired end-to-end interval is available. The
transcriber ordering persists after separation in this configuration. A stronger
transcriber on estimated stems exceeding a weaker model on reference stems is
not evidence of exceeding an oracle bound.

The experiment does not isolate architecture: YourMT3+ and basic-pitch differ in
training data, capacity, decoding, and inference cost. YourMT3+ includes Slakh
training data. Component choices were informed by Slakh test results, so these
are retrospective development observations. A same-checkpoint direct-mixture
YourMT3+ control, with an explicit instrument-to-group mapping, is needed to
measure whether separation adds value.

A separate historical separator comparison holds the transcribers fixed and
reports higher drum SI-SDR for SCNet than HT-Demucs, but lower drum onset F1.
Without paired uncertainty or a controlled mechanism test, this is a diagnostic
observation rather than a statistically established reversal. It motivates
examining downstream notes in addition to listening to separated audio.

\section{Evidence and next steps}
The public repository contains the service, scoring code, experiment ledger, and
aggregate statistics. The per-recording archives supporting this comparison were
unavailable for the present revision and still require release. A committed
summary snapshot regenerates the paper's presentation; a separate checker rejects
missing archives and mismatched track IDs before recomputing stored-score means
and intervals. This does not reproduce inference or establish that stored scores
were computed correctly from MIDI.

The next evaluation should freeze the pipeline on a development split, compare
direct and separated inputs on paired recordings, and measure note offsets,
instrument labels, resource use, and editing effort. A real-audio corpus is needed
for the intended production domain. The demo invites feedback on which visible
errors matter most to musicians and which comparisons would guide component choice.

\section{Acknowledgments and AI use}
We thank the authors and maintainers of the cited datasets, models, and evaluation
software. Claude (Anthropic) generated the original manuscript and assisted
implementation and experiment execution under the author's direction. OpenAI Codex
assisted source checking, evaluation-code auditing, and substantive revision.
This use extends beyond language editing. The named human author is responsible
for the submitted work; neither system is an author. The verification limitations
are stated in the scientific content.
\bibliography{references}
\end{document}
